\chapter{Costochondritis}

\section*{Definition}
Costochondritis is the benign, self-limited inflammation of the costochondral junctions, presenting as anterior chest wall pain reproducible by palpation. Diagnosis is clinical, with focal tenderness of at least one costochondral joint (typically ribs 2--5) in the absence of other pathology.

\section*{Pathophysiology}
The precise mechanism is poorly understood, but costochondritis likely involves microtrauma or repetitive strain to the costochondral cartilage-bone interface, triggering a local inflammatory response with edema and nociceptor sensitization. Unlike Tietze syndrome, costochondritis lacks visible swelling.

\section*{Causes}
\begin{itemize}
  \item \textbf{Idiopathic:} Most common, no identifiable trigger
  \item \textbf{Mechanical:} Heavy lifting, persistent coughing (as in bronchitis or pertussis), strenuous upper body exercise
  \item \textbf{Infectious:} Rare; costochondral infection following thoracic surgery, IV drug use, or in immunocompromised hosts
  \item \textbf{Rheumatologic:} Ankylosing spondylitis, psoriatic arthritis, reactive arthritis, rheumatoid arthritis
  \item \textbf{Other:} Fibromyalgia, chest wall trauma, post-surgical (post-sternotomy)
\end{itemize}

\section*{When to See a Doctor}
See your doctor if:
\begin{itemize}
  \item Chest pain is accompanied by dyspnea, diaphoresis, nausea, or radiation to arm/jaw (rule out cardiac ischemia)
  \item Fever, chills, or erythema/swelling over the costochondral junctions (suggests septic arthritis or Tietze syndrome)
  \item Pain is severe, progressive, or unresponsive to conservative measures
  \item History of cancer with potential for chest wall metastases
  \item Coagulopathy or anticoagulant use with chest wall hematoma suspicion
\end{itemize}

\section*{Related Symptoms}
\begin{itemize}
  \item Anterior chest pain (sharp, stabbing, pressure-like), chest wall tenderness, pain worsened by deep inspiration, cough, or movement
\end{itemize}
\textbf{Important:} Costochondritis is a diagnosis of exclusion; cardiac ischemia must always be considered first in people with cardiac risk factors, regardless of chest wall tenderness.

\section*{Self-Care / Home Management}
\begin{itemize}
  \item NSAIDs (ibuprofen 400--600 mg or naproxen 250--500 mg) taken regularly for 1--2 weeks
  \item Ice packs applied to the affected area for 15--20 minutes several times daily
  \item Avoid activities that exacerbate pain (heavy lifting, push-ups, bench press, twisting motions)
  \item Deep breathing exercises to prevent atelectasis from splinting
  \item Heat therapy after the initial 48 hours to promote blood flow and muscle relaxation
\end{itemize}

\section*{How Your Doctor Will Evaluate You}
\begin{itemize}
  \item Clinical diagnosis; chest wall tenderness reproduced by palpation of costochondral joints is hallmark
  \item ECG and cardiac biomarkers to exclude acute coronary syndrome in at-risk people
  \item Chest X-ray to rule out pneumonia, pneumothorax, or rib fracture
  \item Musculoskeletal ultrasound may reveal costochondral thickening or effusion
  \item Bone scan or MRI for refractory cases to exclude osteomyelitis or tumor
\end{itemize}

\section*{Treatment Options}
\begin{itemize}
  \item NSAIDs (ibuprofen, naproxen, diclofenac gel) as first-line pharmacotherapy for 1--2 weeks
  \item Acetaminophen or topical analgesic creams for people with NSAID contraindications
  \item Physical therapy with stretching exercises (pectoralis major, intercostal muscles) for chronic cases
  \item Corticosteroid and lidocaine injection into the affected costochondral joint for severe, refractory pain
  \item Reassurance and expectant management as most cases resolve spontaneously within weeks to months
\end{itemize}

\clearpage
